\documentclass[a4paper]{article}
\usepackage{amsfonts}
\usepackage{amsmath}
\usepackage{amssymb}
\usepackage{graphicx}     

\begin{document}
  
\noindent \large Auteur: Odia Kibor \\
\noindent \large Studierichting: Informatica\\
\noindent \large Oefeningenreeks 2D nr. 5.9\\

\medskip

\normalsize

$ f(x) = \dfrac{1}{(x - 5)(2x + 1)}$ \\

$\lim\limits_{x \rightarrow 5} f(x) = \dfrac{1}{0} \Rightarrow $ Tekenonderzoek \\

\begin{tabular}{c|cccccc}
$x$     &      &  $\dfrac{-1}{2}$ &       &  $5$  &     \\ \hline
$1$     & $+$  &  $+$  &  $+$     &  $+$  & $+$ \\
$x-5$   & $-$  &  $-$  &  $-$     &  $0$  & $+$ \\
$2x+1$  & $-$  &  $0$  &  $+$     &  $+$  & $+$ \\ \hline
$f(x)$  & $+$  &  $|$  &  $-$     &  $|$  & $+$ \\
\end{tabular}\\

$\Rightarrow \lim\limits_{x \rightarrow 5-} f(x) = -\infty$ \\

$\Rightarrow \lim\limits_{x \rightarrow 5+} f(x) = +\infty$ \\

$\lim\limits_{x \rightarrow -\tfrac{1}{2}} f(x) = \dfrac{1}{0} \Rightarrow $ Tekenonderzoek (zie hierboven)\\

$\Rightarrow \lim\limits_{x \rightarrow -\tfrac{1}{2}-} f(x) = +\infty$ \\

$\Rightarrow \lim\limits_{x \rightarrow -\tfrac{1}{2}+} f(x) = -\infty$ \\










\begin{thebibliography}{999}
%\bibitem{a} Referentie
\end{thebibliography}
\end{document} 


